\section{Conclusion}
In this report, we examined two foundational algorithmic strategies applied to real-world applied computing challenges: greedy interval scheduling for satellite payload management, and divide-and-conquer closest-pair detection for airspace spatial proximity screening. For both problems, we formulated mathematical abstractions, derived asymptotic complexities, proved correctness through formal inductive methods, translated the algorithmic mechanics back into domain concepts, and verified performance via reproducible empirical benchmarking. 

The experimental findings provide empirical evidence consistent with the predicted $\mathcal{O}(n \log n)$ asymptotic scaling across diverse workload sizes, while formal proofs establish theoretical optimality and correctness within the stated operational assumptions. Furthermore, our analysis highlighted essential engineering boundaries---specifically sequence-dependent attitude slewing in satellite operations and 3D kinematic trajectory dynamics in air traffic surveillance---clarifying where these algorithms provide optimal primitives and where higher-layer systems architectures are necessary.
